%=============================================================
%  IEEE Conference style paper  --  METHODOLOGY FOCUS
%
%  Design and Parasitic-Aware Post-Layout Realization of a 6.5 GHz Three-Stage
%  CMOS Ring Oscillator in 65 nm
%
%  COMPILING:
%   * Overleaf / IEEE: keep \documentclass[conference]{IEEEtran} (IEEEtran.cls
%     is preinstalled there; delete the local shim).
%   * Placeholders you must still fill are marked \TODO{...} in red.
%=============================================================
\documentclass[conference]{IEEEtran}

\usepackage{cite}
\usepackage{amsmath,amssymb}
\usepackage{graphicx}
\usepackage{booktabs}
\usepackage{multirow}
\usepackage{array}
\usepackage{textcomp}
\usepackage{xcolor}

\newcommand{\TODO}[1]{\textcolor{red}{\textbf{[TODO: #1]}}}

\begin{document}

\title{Design and Parasitic-Aware Post-Layout Realization\\
of a 6.5~GHz Three-Stage CMOS Ring Oscillator\\
in a 65~nm Technology}

\author{ \IEEEauthorblockN{\TODO{Author Name(s)}}
\IEEEauthorblockA{\TODO{Department, Institution}\\
\TODO{City, Country}\\
Email: \TODO{your.email@domain}} }

\maketitle

%-------------------------------------------------------------
\begin{abstract}
This paper presents the complete custom-design methodology used to realize and
verify a $6.5$~GHz three-stage CMOS ring oscillator in a 65~nm CMOS technology
operating from a $1$~V supply. The emphasis of the work is on the full
physical-design and verification flow that yields a credible, parasitic-aware
oscillation frequency, rather than on idealized schematic estimates. The
oscillator is built from three cascaded static CMOS inverter stages connected in
a feedback loop. The methodology covers manual layout, functional design-rule
checking (DRC), layout-versus-schematic (LVS) verification, and parasitic
resistance--capacitance (RC) extraction in a detailed standard parasitic format
(DSPF). Simulation of the RC-extracted netlist yields an oscillation frequency
of $6.50$~GHz with a power consumption of $286.3~\mu$W from a $1$~V supply,
while occupying a core area of only $73.96~\mu$m$^2$. A periodic steady-state
and periodic-noise analysis of the same RC-extracted netlist gives a phase noise
of $-77.6$~dBc/Hz at a $1$~MHz offset, corresponding to a figure of merit (FoM)
of $-159.3$~dBc/Hz. The flow demonstrates that a parasitic-aware methodology is
essential for an accurate high-frequency ring-oscillator characterization, since
interconnect parasitics substantially lower the oscillation frequency relative
to pre-layout estimates.
\end{abstract}

\begin{IEEEkeywords}
Ring oscillator, CMOS, 65~nm, physical design, DRC, LVS, parasitic extraction,
post-layout simulation, phase noise, DSPF.
\end{IEEEkeywords}

%=============================================================
\section{Introduction}
%=============================================================
Ring oscillators are widely used in modern integrated circuits as the core of
phase-locked loops (PLLs), clock and data recovery circuits, and on-chip clock
generators~\cite{razavi_book}. Compared with $LC$ oscillators, they offer a wide
tuning range, occupy a much smaller die area, and can be implemented in any
standard CMOS process without dedicated passive components~\cite{nugroho2012cco,
pokharel2011quadrature}. Among the possible topologies, the three-stage ring has
been demonstrated successfully in many works~\cite{nugroho2012cco}, which makes
it an attractive and well-understood configuration.

A common shortcoming in reported designs is that performance figures are quoted
from schematic-level simulation, which neglects the parasitics introduced by the
physical layout. In a deep-submicron technology operating at multi-gigahertz
frequencies, interconnect resistance and capacitance significantly affect the
node delays and therefore the oscillation frequency. A meaningful
characterization must therefore be performed on a layout that has been verified
(DRC and LVS) and whose parasitics have been extracted.

This paper focuses on the \emph{methodology} required to obtain such a
parasitic-aware result. The main contributions are: (i)~a complete and
reproducible custom-design flow---layout, functional DRC, LVS, and RC
extraction---for a three-stage 65~nm CMOS ring oscillator; (ii)~a post-layout,
parasitic-aware characterization yielding $6.50$~GHz at $286.3~\mu$W in a
$73.96~\mu$m$^2$ core, including a phase-noise and figure-of-merit
characterization extracted from the same RC netlist; and (iii)~a discussion of
how the parasitics extracted from the layout shape the final oscillation
frequency.

The remainder of the paper is organized as follows. Section~\ref{sec:arch}
describes the circuit architecture. Section~\ref{sec:flow} details the design
and verification methodology. Section~\ref{sec:results} presents the post-layout
simulation results. Section~\ref{sec:conclusion} concludes.

%=============================================================
\section{Circuit Architecture}\label{sec:arch}
%=============================================================
The oscillator is a three-stage inverting ring. Each stage is a static CMOS
inverter consisting of one PMOS pull-up device and one NMOS pull-down device;
the output of the third stage is fed back to the input of the first stage to
close the oscillating loop. The transistor dimensions are summarized in
Table~\ref{tab:sizing}.

\begin{table}[t]
\caption{Transistor Dimensions}
\label{tab:sizing}
\centering
\begin{tabular}{lccc}
\toprule
Device & Type & $W$ & $L$\\
\midrule
NMOS (3 devices) & nch & $2~\mu$m & $60$~nm\\
PMOS (3 devices) & pch & $4~\mu$m & $60$~nm\\
\bottomrule
\end{tabular}
\end{table}

To first order, the period of an $N$-stage ring oscillator is set by the
propagation delay $t_d$ of a single stage,
\begin{equation}
f_{osc} = \frac{1}{2\,N\,t_{d}},\qquad
t_{d}\approx\frac{C_{node}\,\Delta V}{I_{drive}},
\label{eq:freq}
\end{equation}
where $N=3$, $C_{node}$ is the total capacitance loading a node, $\Delta V$ is
the switching voltage, and $I_{drive}$ is the available charging/discharging
current. Equation~\eqref{eq:freq} shows that the oscillation frequency is
inversely proportional to the node capacitance; consequently, the interconnect
parasitics added by the layout directly reduce $f_{osc}$. This dependence
motivates the parasitic-aware methodology described next.

%=============================================================
\section{Design and Verification Methodology}\label{sec:flow}
%=============================================================
The custom flow used in this work is summarized in Fig.~\ref{fig:flow}. After
schematic capture, the circuit was laid out manually, checked against the
functional design rules, verified for schematic equivalence (LVS), and finally
subjected to parasitic RC extraction before post-layout simulation.

\begin{figure}[t]
\centering
\includegraphics[width=0.62\columnwidth]{fig_flow.pdf}
\caption{Custom physical-design and parasitic-aware verification flow used in
this work.}
\label{fig:flow}
\end{figure}

\subsection{Layout}
The ring oscillator was laid out in the target 65~nm CMOS technology. As shown
in Fig.~\ref{fig:layout}, the three PMOS devices are placed in a common $n$-well
along the upper row and share the $V_{DD}$ rail, while the three NMOS devices
occupy the lower row and share the ground rail. The three inter-stage nodes
($V_1$, $V_2$, and the output) are routed in the lower metal layers to keep the
interconnect short. The resulting core occupies an $8.6~\mu\text{m}\times
8.6~\mu\text{m}$ region, i.e. a core area of $73.96~\mu$m$^2$.

\begin{figure}[t]
\centering
\includegraphics[width=0.95\columnwidth]{fig_layout.png}
\caption{Layout of the three-stage CMOS ring oscillator core (8.6~$\mu$m
$\times$ 8.6~$\mu$m). The upper row contains the three PMOS devices and the
lower row the three NMOS devices; the inter-stage nets and the output are routed
between them.}
\label{fig:layout}
\end{figure}

\subsection{Functional Design-Rule Check (DRC)}
DRC was performed to ensure that the drawn geometry complies with the
manufacturing rules of the technology. All \emph{functional} rules, such as
minimum spacing and width, were resolved; for example, a metal-1 spacing
violation (minimum $0.09~\mu$m) was corrected by adjusting the offending
geometry. The remaining DRC items are chip-level seal-ring rules and foundry
density rules (metal, poly, and diffusion density), together with informational
reminders. These pertain to full-chip integration and manufacturing uniformity
rather than to the functional correctness of the oscillator core, and therefore
do not affect the extracted netlist or the simulated behavior.

\subsection{Layout-versus-Schematic (LVS)}
LVS was carried out to confirm that the layout is electrically equivalent to the
schematic. A clean LVS match was obtained, with no shorts, opens, or unmatched
devices. The LVS-clean result is the essential prerequisite for parasitic
extraction and post-layout simulation, since it guarantees that the extracted
netlist faithfully represents the intended topology.

\subsection{Parasitic RC Extraction}
A detailed parasitic extraction was performed on the verified layout, producing
a transistor-level netlist annotated with interconnect resistances and
capacitances in the detailed standard parasitic format (DSPF). The extraction
captures the per-net parasitic capacitances and the metal/poly interconnect
resistances. The extracted per-net capacitances of the key nodes are listed in
Table~\ref{tab:parasitics}; these values quantify the additional loading that
the layout imposes on each oscillating node.

\begin{table}[t]
\caption{Extracted Per-Net Parasitic Capacitance}
\label{tab:parasitics}
\centering
\begin{tabular}{lc}
\toprule
Net & Parasitic capacitance (fF)\\
\midrule
Out   & 4.07\\
$V_{DD}$ & 3.16\\
gnd   & 3.01\\
net3  & 3.54\\
net5  & 3.52\\
\bottomrule
\end{tabular}
\end{table}

The RC-extracted netlist was then simulated under a $1$~V supply using a
transient analysis. To guarantee start-up of the oscillation, an initial
condition was imposed on the output node.

%=============================================================
\section{Post-Layout Simulation Results}\label{sec:results}
%=============================================================
Fig.~\ref{fig:wave_rc} shows the post-layout transient output waveform obtained
from the RC-extracted netlist. The oscillation is stable and spans approximately
the full supply range. Two adjacent peaks were measured at $t_1 = 137.77$~ps and
$t_2 = 291.61$~ps, giving a period
\begin{equation}
T = t_2 - t_1 = 153.84~\text{ps},
\end{equation}
and an oscillation frequency
\begin{equation}
f_{osc} = \frac{1}{T} \approx 6.50~\text{GHz}.
\end{equation}

\begin{figure}[t]
\centering
\includegraphics[width=0.95\columnwidth]{fig_wave_rc.png}
\caption{Post-layout (RC-extracted) transient output waveform. The two markers
at $137.77$~ps and $291.61$~ps yield a period of $153.84$~ps, i.e. an
oscillation frequency of $6.50$~GHz.}
\label{fig:wave_rc}
\end{figure}

The average supply current was obtained from the transient supply current,
$|I_{avg}| = 286.3~\mu$A, giving a power consumption of
\begin{equation}
P = V_{DD}\,|I_{avg}| = 1~\text{V}\times 286.3~\mu\text{A}
= 286.3~\mu\text{W}.
\end{equation}
The output swing extends over approximately the full supply range (about
$-0.05$~V to $1.04$~V), and the per-stage delay implied by \eqref{eq:freq} is
$t_d = T/(2N) = 25.6$~ps. The complete set of measured performance metrics is
summarized in Table~\ref{tab:results}.

\begin{table}[t]
\caption{Summary of Post-Layout (RC) Performance}
\label{tab:results}
\centering
\begin{tabular}{lc}
\toprule
Parameter & Value\\
\midrule
Technology              & 65~nm CMOS\\
Supply voltage $V_{DD}$ & $1.0$~V\\
Topology                & 3-stage CMOS ring\\
Oscillation frequency   & $6.50$~GHz\\
Period                  & $153.84$~ps\\
Stage delay             & $25.6$~ps\\
Power consumption       & $286.3~\mu$W\\
Output swing            & $\sim\!1.0$~V\\
Phase noise @ 1~MHz     & $-77.6$~dBc/Hz\\
FoM                     & $-159.3$~dBc/Hz\\
Core area               & $73.96~\mu$m$^2$\\
\bottomrule
\end{tabular}
\end{table}

\subsection{Phase Noise and Figure of Merit}
To characterize the spectral purity of the oscillator, a periodic steady-state
(PSS) analysis followed by a periodic-noise (pnoise) analysis was performed
directly on the RC-extracted netlist. The resulting single-sideband phase-noise
characteristic is shown in Fig.~\ref{fig:pnoise}. The curve exhibits the
expected $1/f^2$ slope across the offset range, confirming a physically
consistent result. At a $1$~MHz offset from the $6.5$~GHz carrier, the phase
noise is $-77.6$~dBc/Hz, decreasing to $-124.3$~dBc/Hz at a $100$~MHz offset.

\begin{figure}[t]
\centering
\includegraphics[width=0.95\columnwidth]{fig_pnoise.png}
\caption{Post-layout (RC-extracted) single-sideband phase noise versus offset
frequency. The marker indicates $-77.6$~dBc/Hz at a $1$~MHz offset.}
\label{fig:pnoise}
\end{figure}

A widely used figure of merit (FoM) for oscillators normalizes the phase noise
to the oscillation frequency and the power consumption~\cite{nugroho2012cco}:
\begin{equation}
\mathrm{FoM} = \mathcal{L}(\Delta f)
- 20\log_{10}\!\left(\frac{f_{osc}}{\Delta f}\right)
+ 10\log_{10}\!\left(\frac{P_{dc}}{1~\text{mW}}\right),
\label{eq:fom}
\end{equation}
where $\mathcal{L}(\Delta f)$ is the phase noise at offset $\Delta f$. Using
$\mathcal{L}(1~\text{MHz}) = -77.6$~dBc/Hz, $f_{osc} = 6.5$~GHz, $\Delta f =
1$~MHz, and $P_{dc} = 286.3~\mu$W, the evaluation of \eqref{eq:fom} gives
\begin{equation}
\mathrm{FoM} = -77.6 - 76.26 - 5.43 \approx -159.3~\text{dBc/Hz}.
\end{equation}
This value is competitive with previously reported CMOS ring
oscillators~\cite{nugroho2012cco}, indicating that the compact, parasitic-aware
design achieves a favorable balance between oscillation frequency, power
consumption, and spectral purity.

\subsection{Discussion}
The post-layout frequency of $6.50$~GHz reflects the additional capacitive and
resistive loading introduced by the interconnect, as quantified in
Table~\ref{tab:parasitics}. In accordance with \eqref{eq:freq}, the increased
node capacitance raises the stage delay and therefore lowers the oscillation
frequency relative to a pre-layout estimate. This confirms that, at
multi-gigahertz frequencies in a 65~nm process, parasitic-aware post-layout
simulation is necessary to obtain a realistic figure; a schematic-only estimate
would significantly overstate the achievable frequency.

Despite the compact $73.96~\mu$m$^2$ core, the oscillator delivers a $6.50$~GHz
output while consuming only $286.3~\mu$W, demonstrating that the proposed
methodology yields a verified, area- and power-efficient high-frequency
oscillator suitable for integration in PLLs and on-chip clock generators.
\TODO{Optionally add a comparison table against prior art if your supervisor
wants it; data for several references are available in \cite{nugroho2012cco}.}

%=============================================================
\section{Conclusion}\label{sec:conclusion}
%=============================================================
A three-stage CMOS ring oscillator was designed and realized through a complete,
parasitic-aware physical-design methodology in a 65~nm CMOS technology. The
flow---manual layout, functional DRC, LVS, and RC extraction---produced a
verified layout from which a realistic post-layout characterization was
obtained. The RC-extracted design oscillates at $6.50$~GHz while consuming
$286.3~\mu$W from a $1$~V supply and occupying only $73.96~\mu$m$^2$. The
results highlight the importance of parasitic-aware post-layout simulation for
accurate high-frequency ring-oscillator design.

%=============================================================
\section*{Acknowledgment}
%=============================================================
\TODO{Optional: acknowledgments (supervisor, laboratory, institution, funding).}

%=============================================================
\begin{thebibliography}{1}

\bibitem{nugroho2012cco}
P.~Nugroho, R.~K. Pokharel, H.~Kanaya, and K.~Yoshida, ``A 5.9~GHz low power and
wide tuning range CMOS current-controlled ring oscillator,'' \emph{Int.\ J.\
Electr.\ Comput.\ Eng.\ (IJECE)}, vol.~2, no.~3, pp.~293--300, Jun. 2012.

\bibitem{razavi_book}
B.~Razavi, \emph{Design of Analog CMOS Integrated Circuits}. Singapore:
McGraw-Hill, 2001.

\bibitem{pokharel2011quadrature}
R.~K. Pokharel \emph{et al.}, ``A wide tuning range CMOS quadrature ring
oscillator with improved FoM for inductorless reconfigurable PLL,'' \emph{IEICE
Trans. Electron.}, vol.~E94-C, no.~10, pp.~1524--1532, 2011.

\bibitem{razavi_phasenoise}
B.~Razavi, ``A study of phase noise in CMOS oscillators,'' \emph{IEEE J.
Solid-State Circuits}, vol.~31, no.~3, pp.~331--343, Mar. 1996.

\end{thebibliography}

\end{document}